\documentclass[10pt,a4paper]{amsart}
\usepackage[margin=19mm]{geometry}
\usepackage{amsmath,amssymb,mathtools}
\usepackage[scr=rsfs,cal=euler]{mathalfa}
\usepackage{xfrac}
\usepackage[unicode,hidelinks,bookmarks=false]{hyperref}
\newcommand{\ie}{i.e.\ }
\newcommand{\cf}{cf.\ }
\newcommand{\resp}{respectively\ }
\newcommand{\Subsection}{\subsection}
\providecommand{\nobreakdash}{\nobreak\hbox{-}}
\setlength{\emergencystretch}{2em}
\multlinegap0pt
\usepackage{amssymb}
\usepackage{mathtools}
\usepackage{xfrac}
\usepackage{float}
\usepackage{xcolor}
\usepackage[shortlabels]{enumitem}
\usepackage{tikz}
\newtheorem{proposition}{Proposition}
\newtheorem{lemma}{Lemma}
\newcommand{\F}{\mathbb{F}}
\newcommand{\Z}{\mathbb{Z}}
\newcommand\sm[1]{\begin{bsmallmatrix}#1\end{bsmallmatrix}}
\title{Number of $K$-rational points with given $j$-invariant on modular curves}
\author{Ivan Novak}
\date{Source: arXiv:2512.24817v2 (CC BY 4.0)}
\begin{document}
\maketitle

\noindent\emph{Source and licence.} Number of $K$-rational points with given $j$-invariant on modular curves. Version arXiv:2512.24817v2 (CC BY 4.0), distributed by arXiv under the Creative Commons Attribution 4.0 International licence, http://creativecommons.org/licenses/by/4.0/, as recorded in the arXiv licence metadata for this exact version and shown on the abstract page https://arxiv.org/abs/2512.24817v2. Full licence notice: \texttt{licenses/arxiv-2512.24817-CC-BY-4.0.txt}.\par\smallskip

\noindent\textit{Editorial scope.} Counted unit: the article's proposition proposition (source0.tex lines 716--742). The article's complete original proof of it is delivered with it (source0.tex lines 744--908, one \texttt{proof} environment of 165 lines). No mathematics is written, repaired or re-derived: the statement and the proof are the article's own lines, in the article's own notation, and no retained line is substituted or re-flowed.  No further source-local context is needed: every cross-reference of every retained line resolves inside the two delivered spans.  Nothing was omitted from any retained span, so the package carries no omission record.  No retained line contains a citation, so the wrapper carries no bibliography and no bibliography entry.  No retained line contains a figure environment, an image-inclusion command or drawing code, so no image is delivered and the package declares no asset dependency.  Editorial formatting only, in addition: an AMS \texttt{amsart} preamble that restates the article's own declarations and macros used by the retained lines, namely the environments \texttt{proposition}, \texttt{lemma} on separate counters, together with the standard packages those lines need.  The retained environments print in this document as Proposition 1, Lemma 1 in order of appearance, whereas the article prints the same items with the numbering of its own full document.  The complete native source of the article as distributed in the arXiv e-print for this exact version is bundled unchanged as \texttt{sources/191932/source0.tex}.
\begin{proposition}
     Let $p$ be an odd prime and $k$ a positive integer.  Let $$
\begin{aligned}
A'_k(p)&=\{0\} \cup \left\{ \frac{p^{2k}-p^{2k-1}}{2} \right\}\cup \{p^{2j} \mid j\in \{0,1,\ldots,k-1\}\},\\
B'_k(p)&=\left\{\frac{p^k+p^{k-1}}{2}+1\right\},\\
C'_k(p)&=\begin{cases}
\{2\}, & \text{if } p\equiv 1\pmod{4},\\[4pt]
\{3\}, & \text{if } p\equiv 3\pmod{4},
\end{cases}\\
D'_k(p)&=\left\{1,\ \frac{p^k+p^{k-1}}{2},\ \frac{p^k-p^{k-1}}{2}\right\},\\
E'_k(p)&=\{p^j \mid 0<j<k\},\\
F'_k(p)&=\begin{cases}
\{p^{\lfloor k/2 \rfloor}+1\}\ \cup\ \{\,2p^t+1 \mid 0<2t<k\,\}, & \text{if } p\equiv \pm 3 \pmod 8,\\[4pt]
\emptyset, & \text{if } p\equiv \pm 1 \pmod 8,
\end{cases}\\
G'_k(p)&=\begin{cases}
\{1,2,3\}, & \text{if } p>3,\\[4pt]
\{1,3\}, & \text{if } p=3,
\end{cases}\\
H'_k(p)&=\begin{cases}
\{1,2\}, & \text{if } p\equiv 1\pmod{4},\\[4pt]
\{1\}, & \text{if } p\equiv 3\pmod{4}.
\end{cases}
\end{aligned}
$$ Then 
    $$S(N_{\mathrm{ns}}(p^k))=\bigcup_{L=A}^H L'_k(p).$$
\end{proposition}

\begin{proof}
    Let $H=N_{\mathrm{ns}}(p^k)$. Let $R$ be a subgroup of $H$. The list of coset representatives for $N_{\mathrm{ns}}(p^k)$ is the same as the one for $C_{\mathrm{ns}}(p^k)$, except $\sm{1 & v \\ 0 & z}$ and $\sm{1 & v \\ 0 & -z}$ yield the same coset. We will find the number of matrices $g=\sm{1 & v \\ 0 & z}$ satisfying $gRg^{-1}\subseteq H$, and then divide this number by $2$.
    
   We have already calculated the form of conjugates in $C_{\mathrm{ns}}(p^k)$ by $\sm{1 & v \\ 0 & z}$ as 
    $$\begin{pmatrix} az+bvz & b(\epsilon - v^2) \\
bz^2 & az-bvz\end{pmatrix}.$$
Similarly, calculating $\sm{1 & v \\0 & z}\sm{a & \epsilon b \\ -b  & -a}\sm{1 & v \\ 0 & z}^{-1}$ yields 
  $$\begin{pmatrix} az-bvz & bv^2-2av+\epsilon b \\
-bz^2 & bvz-az\end{pmatrix}.$$

\noindent \textbf{Case I.} $R$ is contained in $C_{\mathrm{ns}}(p^k)$. 

\noindent \textbf{Case I.a.} For every solution $g$, one has $gRg^{-1}\subseteq C_{\mathrm{ns}}(p^k)$. This case is equivalent to the case from part (c), and we obtain the same set of numbers, except we divide them by $2$ because the index is twice smaller.

\noindent \textbf{Case I.b.} There exists a solution $g$ with $gRg^{-1}\not \subseteq C_{\mathrm{ns}}(p^k)$. This means that $R$ contains a matrix $m=\sm{a & b \\ \epsilon b & a}$ such that  $$\begin{pmatrix} az+bvz & b(\epsilon - v^2) \\
bz^2 & az-bvz\end{pmatrix}.$$ belongs to $H\setminus C_{\mathrm{ns}}(p^k)$. This is equivalent to 
\[
\begin{cases}
2az=0, \\[4pt]
b(\epsilon z^2+\epsilon)=bv^2.
\end{cases}
\]
Since $2z$ is invertible, it follows that $a=0$. Hence, $b$ is invertible and we get the equation \begin{equation} \label{bla1}\epsilon z^2+\epsilon=v^2.\end{equation}

Furthermore, if $g$ is a solution such that $gmg^{-1}$ lies in $C_{\mathrm{ns}}(p^k)$, it follows from the proof of part (c) that $g\in N_{\mathrm{ns}}(p^k)$ (note that $\beta=0$ in that case). Thus, there is exactly one class of solutions $g$ such that $gmg^{-1} \in C_{\mathrm{ns}}(p^k)$. Any other solution satisfies \eqref{bla1}.

Now suppose that $R$ contains a nonscalar matrix $\sm{a & \epsilon b \\b & a}$ with $a\neq 0$. This means that the conjugate of this matrix by $\sm{1 & v \\ 0 & z}$ lies in $C_{\mathrm{ns}}(p^k)$. Part (c) then implies that $v\equiv 0 \pmod p$ and $z^2\equiv 1 \pmod p$. However, then \eqref{bla1} implies $2\epsilon \equiv 0 \pmod p$, a contradiction. Thus, it suffices to count the number of solutions of the equation \eqref{bla1} for which $z$ is invertible. However, note that $z$ is always invertible, since otherwise we would have $\epsilon \equiv v^2 \pmod p$.

It follows from a standard character calculation argument that the number of solutions $(v,z)$ modulo $p$ to  \eqref{bla1} is $p-\chi_p(\epsilon)=p+1$, where $\chi_p$ is the quadratic character mod $p$. To go from $p$ to $p^k$, note that $z$ is always invertible, so Hensel's lemma implies that any solution in $\Z/p\Z$ has exactly $p^{k-1}$ lifts to $\Z/p^k\Z$. In total, after dividing by $2$ to account for the smaller index, we obtain $\frac{p^k+p^{k-1}}{2}$ solutions. However, we must also add the number of matrices $m=\sm{1 & v \\ 0 & z}$ such that $m\sm{0 & \epsilon b \\\ -b & 0}m^{-1}$ lies in $C_{\mathrm{ns}}(p^k)$. This is equivalent to $v=0$ and $z^2=1$, and since $z$ and $-z$ define the same coset, the total number of solutions is $\frac{p^k+p^{k-1}}{2}+1$.

This number is achieved by taking $R$ to be the group generated by $\sm{0 & \epsilon  \\ 1 & 0}.$

\noindent \textbf{Case II.} $R$ is not contained in $C_{\mathrm{ns}}(p^k)$. Note that $R_1:=R\cap C_{\mathrm{ns}}(p^k)$ is an index $2$ subgroup of $R$. We know that either any solution $g$ satisfies $gR_1g^{-1}\subseteq C_{\mathrm{ns}}(p^k)$, or any matrix $\sm{a & \epsilon b \\b & a}$ in $R_1$ satisfies either $a=0$ or $b=0$. 

Let $R_2=R\setminus R_1$. Then $R_2=rR_1$ for some $r=\sm{a_0 & \epsilon b_0 \\ -b_0 & -a_0}\in R_2$.  

\noindent \textbf{Case II.a.} There exists a solution $g$ such that $grg^{-1}\in C_{\mathrm{ns}}(p^k)$. We claim that $R_1$ consists only of matrices $\sm{a & \epsilon b \\ b & a}$ such that $a=0$ or $b=0$.


 Since $grg^{-1} \in C_{\mathrm{ns}}(p^k)$, we have $$\begin{cases}
    a_0z=b_0vz, \\
    -\epsilon b_0z^2=b_0v^2-2a_0v+\epsilon b_0.
\end{cases}    
$$
Since $z$ is invertible, this simplifies to $$\begin{cases}
    a_0=b_0v, \\
    -\epsilon b_0z^2=-b_0v^2+\epsilon b_0.
\end{cases}$$
Note that $b_0$ is invertible, since otherwise $a_0$ would not be. Thus, $-\epsilon z^2+v^2=\epsilon.$
Consider a matrix $m=\sm{a & \epsilon b \\ b & a}$ in $R_1$. If $gmg^{-1}\not \in C_{\mathrm{ns}}(p^k)$, then $a=0$. Otherwise, if $b\neq 0$, we can use the same argument as in part (c), to obtain $v\equiv 0 \pmod p$ and $z^2\equiv 1 \pmod p$. However, this implies $-\epsilon \equiv \epsilon \pmod p$, a contradiction.   

Thus, each matrix $\sm{a &  \epsilon b \\ b & a}$ in $R_1$ satisfies either $a=0$ or $b=0$.

\noindent \textbf{Case II.a.(i)} Each matrix in $R_1$ is a scalar matrix. This means that a solution is any matrix $g$ such that $grg^{-1} \in C_{\mathrm{ns}}(p^k)$ or $grg^{-1} \in H\setminus C_{\mathrm{ns}}(p^k)$. Thus, we need to find the sum of the numbers of solutions to the system $$\begin{cases}
a_0=b_0v, \\
v^2=\epsilon z^2+\epsilon,
\end{cases}$$
and add to that number the number of solutions of the equation\begin{equation} \label{bla2}
    \epsilon b_0 z^2=b_0 v^2-2a_0v+\epsilon b_0,
\end{equation}
with $a_0$ and $b_0$ fixed and at least one of them invertible, and with the assumption that the first system has at least one solution. Note that in fact the first system has at most one solution, since $v$ is uniquely determined from the first equation, and $z^2$ is uniquely determined from the second equation. Also, this equation does have a solution if and only if $$\frac{a_0^2}{\epsilon b_0^2}-1$$ is a square modulo $p$. Note that this is always possible to achieve, as there always exist elements $x$ and $x+1$ in $(\Z/p\Z)^\times$ such that $x$ is a residue and $x+1$ a nonresidue.

It remains to find the number of solutions to \eqref{bla2} for which $z$ is invertible. We can divide the equation by $b_0$ since it is also invertible. We get $$\epsilon z^2=v^2-2\frac{a_0}{b_0}v+\epsilon.$$
Again, using standard character counting argument, it follows that there are $p+1$ solutions modulo $p$ to this equation. Note that there are no solutions with $z=0$, since we would have $$\left(v-\frac{a_0}{b_0} \right)^2=\epsilon\left(\frac{a_0^2}{\epsilon b_0^2}-1\right),$$ and the right hand side is not a square mod $p$. Thus, we get $p+1$ solutions modulo $p$ in this case. There are $p^{k-1}$ lifts of each solution, so we obtain $p^k+p^{k-1}$ solutions.  Dividing by $2$ to account for $z$ and $-z$ being the same, and adding $1$ for the solution of the system, we again obtain the number $\frac{p^k+p^{k-1}}{2}+1$.

\noindent \textbf{Case II.a.(ii)} There exists a nonscalar matrix in $R_1$.

As in the previous case, we conclude that there is exactly one solution to $$\begin{cases}
a_0=b_0v, \\
v^2=\epsilon z^2+\epsilon,
\end{cases}.$$

Consider a nonscalar matrix $\sm{0 & \epsilon b \\ b & 0} \in R_1$. Note that if $\sm{1 & v \\0 & z}$ is the solution to the system above, we have $\sm{1 & v \\0 & z }\sm{0 & \epsilon b \\ b & 0}\sm{1 & v \\ 0 & z}^{-1}\in H\setminus C_{\mathrm{ns}}(p^k)$, so the solution to the system remains valid. Also, there is always the trivial solution $v=0$, $z^2=1$. 

Suppose that there is some other nontrivial solution $(v_0, z_0)$. Then $\epsilon z_0^2+\epsilon= v_0^2$ by looking at the conjugate of $\sm{0 & \epsilon b \\ b & 0 }$. On the other hand, looking at the conjugate of $r=\sm{a_0 & \epsilon b_0 \\ -b_0 & -a_0}$, it follows that $$\epsilon b_0 z_0^2=b_0 v_0^2-2a_0v_0+\epsilon b_0.$$ We again divide by $b_0$, and then eliminate $\epsilon z^2$ from the equation to obtain $$v_0^2-\epsilon=v_0^2-2\frac{a_0}{b_0}v_0+\epsilon,$$ from where it follows that $$v_0=\frac{\epsilon b_0}{a_0}.$$
Plugging this back in the expression $\epsilon z_0^2+\epsilon=v_0^2$, it follows that $$\epsilon z_0^2+\epsilon=\frac{b_0^2\epsilon^2}{a_0^2},$$ and finally $$z_0^2=\frac{\epsilon b_0^2}{a_0^2}-1.$$
This equation has a solution if and only if $\frac{\epsilon b_0^2}{a_0^2}-1$ is a square, which is equivalent with $\epsilon b_0^2-a_0^2$ being a square. Since $\frac{a_0^2}{\epsilon b_0^2}-1$ is a square, it follows that $a_0^2-\epsilon b_0^2$ is not a square. Thus, it depends on whether $-1$ is a square or not. 

If $-1$ is a square, then $\epsilon b_0^2-a_0^2$ is not a square and we obtain $2$ solutions in total. 

If $-1$ is not a square, then $\epsilon b_0^2-a_0^2$ is a square and we obtain an additional solution, for a total of $3$ solutions.

\noindent \textbf{Case II.b.} All solutions $g$ satisfy $grg^{-1} \in H\setminus C_{\mathrm{ns}}(p^k)$. In this case, any solution satisfies the equation \eqref{bla2}, which we write again for clarity reasons: $$\epsilon b_0 z^2=b_0 v^2-2a_0v+\epsilon b_0.$$

\noindent \textbf{Case II.b.(i)} All solutions $g$ also satisfy $gR_1g^{-1}\subseteq C_{\mathrm{ns}}(p^k)$. 

Let $\beta=\min_{g \in R_1}v_p(b)$ as before. If $\beta=0$, we obtain only the trivial solution, as before. Otherwise, we need to count the number of solutions to \eqref{bla2} which also satisfy $p^{k-\beta} \mid z^2-1$ and $p^{k-\beta} \mid v$.

We start by solving the special case $\beta=k$. In that case, we only need to count the number of spolutions to \eqref{bla2}. 

If $b_0$ is not invertible, then $a_0$ is invertible. Note that then $v=0$ is a solution modulo $p$, and Hensel's lemma implies that this solution has a unique lift, for any value of $z^2$. There are $\frac{p^{k-1}(p-1)}{2}$ choices for $z^2$, so this is the number of solutions.

If $b_0$ is invertible, then using the same arguments as before, there are $p+1$ solutions $(z,w)$ mod $p$. Of these solutions, there are either $0$ or $2$ solutions with $z=0$ which we need to exclude, depending on whether $4a_0^2-4\epsilon b_0$ is a square or not. If $4a_0^2-4\epsilon b_0$ is a square, then we obtain $\frac{(p-1)p^{k-1}}{2}$ solutions as before. If it is not a square, we obtain $\frac{(p+1)p^{k-1}}{2}$ solutions.

Now suppose $0<\beta<k$. Note that there is then only one solution mod $p$, which is $z^2=1$, $v=0$. 

If $b_0$ is not invertible, then again by Hensel's lemma, for any choice of $z$ such that $z^2-1$ is divisible by $p^{k-\beta}$, there is a unique lift of the solution $v=0$ to a solution modulo $p^k$. It is easy to see from the equation that this lift satisfies $p^{k-\beta} \mid v$. Thus, in this case we get $p^{\beta}$ solutions, one for each choice of $z^2$.

If $b_0$ is invertible, then for any choice of $v$ such that $p^{k-\beta} \mid v$, there is a unique lift of the solution $z^2=1$ to a solution modulo $p^k$. Again, it is easy to see from the equation that this lift satisfies $p^{k-\beta} \mid z^2-1$. Again, we obtain $p^\beta$ solutions.


\noindent \textbf{Case II.b.(ii)} Every matrix $\sm{a & \epsilon b \\ b & a}\in R_1$ satisfies $a=0$ or $b=0$, and there exists a matrix with $a=0$. In this case, any nontrivial solution $g$ satisfies $g \sm{0 & \epsilon b \\ b & 0}g^{-1}\in H\setminus C_{\mathrm{ns}}(p^k)$, which gives us the equality
$\epsilon z^2+\epsilon =v^2$. This allows us to eliminate $z^2$ from \eqref{bla2}, and obtain $$-2\epsilon b_0=b_0 v^2-2a_0v.$$

If $b_0$ is invertible, the equation can be rewritten as $$v^2-2\frac{a_0}{b_0}v+2\epsilon=0.$$

This equation has a double root modulo $p$ if and only if $\frac{a_0^2}{b_0^2}\equiv 2\epsilon \pmod p$. If $2$ is a square modulo $p$, this is impossible and we only get the trivial solution. 

If $2$ is not a square modulo $p$, then this is possible. In that case, $v\equiv \frac{a_0}{b_0} \pmod p$ and $\frac{a_0^2}{b_0^2}\equiv 2\epsilon \pmod p$. Note that this implies $z^2\equiv 1 \pmod p$, so for any value of $v$ that satisfies the conditions, there is a unique suitable value of $z$.

It suffices to count the number of solutions to $$v^2-2\frac{a_0}{b_0}v+\frac{a_0^2}{b_0^2}-px=0,$$ where $p x=\frac{a_0^2}{b_0^2}-2\epsilon$. If $s=v-\frac{a_0}{b_0}$, this is equivalent with $$s^2=px.$$
There may be $0$ solutions to this equation. Suppose that there is at least $1$ solution. 

If $px=0$, there are $p^{\lfloor k/2 \rfloor}$ solutions. Adding $1$ from the trivial solution, we obtain $p^{\lfloor k/2 \rfloor}+1$. 

Otherwise, let $px=p^{2t}c^2$ where $2t<k$ and $c$ invertible. Then $v_p(s)=t$, and $s=p^t r$ for some invertible $r$, and $r^2\equiv c^2 \pmod p^{k-2t}$. There are $2$ choices for $r \pmod {p^{k-2t}}$, hence $2p^{t}$ choices mod $p^{k-t}$. In total, there are $2p^{t}$ choices for $v$. Adding $1$ from the trivial solution, we obtain $2p^t+1$.

If the equation has only single roots, then its discriminant needs to be a square modulo $p$ to have solutions. That is, $\frac{a_0^2}{b_0^2}-2\epsilon$ needs to be a square. As before, this can be achieved for some $a_0$ and $b_0$.  In that case, there are exactly $2$ solutions $v_1$ and $v_2$ by Hensel's lemma. It remains to check whether $z^2=\frac{v_i^2}{\epsilon}-1$ is a square. Depending on whether none, one or both of them is a square, we obtain $1$, $2$ or $3$ solutions respectively. We claim that we can always choose $\frac{a_0}{b_0}$ for each of the possibilities when $p$ is large enough. 

Let $2\frac{a_0}{b_0}=c$. Then $c^2-8\epsilon$ is a square. Let $d^2=c^2-8\epsilon$. Then $\frac{v_1^2}{\epsilon}-1=\frac{(d+c)^2}{4\epsilon}-1$ and $\frac{v_2^2}{\epsilon}-1=\frac{(d-c)^2}{4\epsilon}-1.$

Multiplying everything by suitable squares and setting $c-d=A$ and $c+d=\frac{8\epsilon }{A}$, we obtain the following two expressions: $16\epsilon-A^2$ and $\epsilon A^2-4\epsilon^2$. 

\begin{lemma}
    Let $p>11$ be a prime, and let $f(x)$ and $g(x)$ be irreducible quadratic polynomials in $\F_p[x]$ with no common factors. Then there is some $A\in \F_p$ such that $f(A)$ and $g(A)$ are both squares in $\F_p^\times$.
\end{lemma}
\begin{proof}
Write \(\chi\) for the quadratic character modulo \(p\).

For signs \(s,t\in\{+1,-1\}\) let
\[
N_{s,t}=\#\{A\in\F_p:\ \chi(f(A))=s,\ \chi(g(A))=t\}.
\]
Using the standard identity \(\mathbf{1}_{\chi(f(x))=s}=(1+s\chi(f(x)))/2\) (and similarly for \(\chi_g\)) we obtain
\[
N_{s,t}=\sum_{A\in\F_p}\frac{1+s\chi(f(A))}{2}\cdot\frac{1+t\chi(g(A))}{2}
=\frac{1}{4}\Big(p+s\Sigma_f+t\Sigma_g+st\Sigma_{fg}\Big),
\]
where
\[
\Sigma_f=\sum_{A\in\F_p}\chi(f(A)),\qquad
\Sigma_g=\sum_{A\in\F_p}\chi(g(A)),\qquad
\Sigma_{fg}=\sum_{A\in\F_p}\chi(f(A)g(A)).
\]


For an irreducible quadratic polynomial \(q(x)=ax^2+bx+c\) one has \(\sum_{x\in\F_p}\chi(q(x))=-\chi(a).\) In particular, \(|\Sigma_f|=|\Sigma_g|=1\).  
The product \(f(x)g(x)\) is a quartic polynomial with no repeated roots, so the Weil bound gives us
\[
|\Sigma_{fg}|\leq  3\sqrt p.
\]

Combining these estimates we obtain, for every sign pair \((s,t)\),
\[
N_{s,t}\ge \frac{1}{4}\Big(p - \big(|\Sigma_f|+|\Sigma_g|+|\Sigma_{fg}|\big)\Big)
\ge \frac{1}{4}\big(p - 2 - 3\sqrt p\big).
\]
The right-hand side is positive whenever \(p-2-3\sqrt p>0\), which holds for every prime \(p\ge 13\). Hence for \(p\ge 13\) each \(N_{s,t}\) is positive. In particular, there exists \(A\) with \(\chi_f(A)=\chi_g(A)=+1\), i.e.\ both \(f(A)\) and \(g(A)\) are nonzero squares.
\end{proof}

Applying this for polynomials $16\epsilon-x^2$ and $\epsilon x^2-4\epsilon^2$, we obtain that each combination of squares and non-squares occurs. It remains to check primes up to $11$. A direct check shows that for $p=3$ either both polynomials are squares or non-squares, while for $p=5,7,11$ there can be $0,1$ or $2$ squares. In conclusion, adding the trivial solution, we obtain $1$, $2$ or $3$ solutions, unless $p=3$ in which case we obtain $1$ or $3$ solutions.


Now suppose that $b_0$ is not invertible. Then $a_0$ is invertible and $v$ is divisible by $p$. By Hensel's lemma, since the derivative of $b_0v^2-2a_0v+2\epsilon b_0$ at $v=0$ is nonzero, it follows that there is a unique solution $v$, and $v\equiv 0 \pmod p$. It follows that $\epsilon z^2+\epsilon\equiv 0 \pmod p$, so there is $1$ nontrivial solution if $-1$ is a square, and $0$ nontrivial solutions otherwise. In any case, we obtain either $1$ or $2$.
\end{proof}

\end{document}
